\Needspace{18\baselineskip}
\subsection{One-way hypernets and the term-boundary restriction}
\label{sec:one-way}
This milestone proves the single-hole case of Appendix~C, Lemma~26(2), for
actual shallow substitution and the concrete rooted-state predicate. The
proof follows the paper's reason for introducing one-way hypernets: an
operation path from the global input cannot pass through such a filling and
emerge at a term output. We make the path-extraction argument explicit and
allow the replaced edge to have a self-loop.

\begin{lemma}[Sources of operation paths have term type]
Every source vertex of an actual operation edge has type $\star$. If a
nonempty reverse operation trace ends at an input, that input has type
$\star$.
\end{lemma}
\begin{proof}
The signature of every primitive operation has exactly one source, of type
$\star$. Apply the source-typing equation to any member of its source list.
For the second assertion, induct on the trace's edge list. If its tail is
empty, its endpoint is the source of the first operation and the first
assertion applies. Otherwise apply the induction hypothesis to the tail.
Boxes, variables, focuses, and holes are excluded by the operation predicate.
\end{proof}
\lean{Graph/OneWay}{10}{22}{Linear.operation_source_term}
\lean{Graph/OneWay}{26}{36}{Linear.operationTrace_input_term}

\begin{definition}[One-way hypernet]\label{term:one-way}
A hypernet $H:I\to O$ is \emph{one-way} if, for every term input $i$ and term
output $j$, there is no operation path from $i$ to $j$. Paths to thunk outputs
are deliberately not excluded. This is the restriction in the paper's
Section~6.1, expressed using our existing operation-path datatype.
\end{definition}
\lean{Graph/OneWay}{40}{43}{Linear.OneWay}

\begin{theorem}[Reachability characterization]
A hypernet is one-way if and only if none of its term outputs is
operation-reachable from an input.
\end{theorem}
\begin{proof}
A reverse trace from an output to an input cannot be empty: the boundary
constructors are disjoint. Its input therefore has term type by the preceding
lemma. Reverse its recorded edge list to obtain the forbidden forward
operation path. Conversely any forbidden forward path supplies precisely a
reverse trace witnessing operation reachability of its term output.
\end{proof}
\lean{Graph/OneWay}{46}{54}{Linear.oneWay_iff}

One-wayness is invariant under recursive graph isomorphism. Such an
isomorphism fixes the ordered boundary and preserves and reflects operation
reachability, so the characterization applies in both directions. A graph
with no outputs is one-way vacuously, since its output index type is empty.
\lean{Graph/OneWay}{66}{73}{Linear.DeepIso.oneWay_iff}
\lean{Graph/OneWay}{76}{78}{Linear.oneWay_no_outputs}

\Needspace{18\baselineskip}
\subsection{Extracting an operation path from a filling}
Fix a host $C$, a selected edge $h$, and a filling $H$ with the interface of
$h$. Write $P=C[H]$, $\iota_C$ for the retained-host vertex map, and
$\iota_H$ for the filling map. These are the existing maps of shallow
replacement; $\iota_H$ need not be injective when $h$ has a self-loop.

The elementary decomposition used in the induction says that a nonempty
reverse operation trace with edge list $e::E$, from $v$ to $z$, consists
exactly of an intermediate vertex $w$, an operation label on $e$, incidences
$v\in\operatorname{target}(e)$ and $w\in\operatorname{source}(e)$, and a
trace with edge list $E$ from $w$ to $z$. Both directions follow directly
from the trace constructor.
\lean{Graph/OneWay}{57}{63}{Linear.operationTrace_cons_iff}

\begin{lemma}[An inserted non-input stays a non-input]
If a vertex $v$ of $H$ is not an input, its image $\iota_H(v)$ is not an
input of $P$.
\end{lemma}
\begin{proof}
Target coverage in $H$ supplies an edge targeting $v$. Its inserted copy
targets $\iota_H(v)$ in $P$. Inputs of a well-formed graph have no incoming
edge, which proves the claim, including when $v$ is a filling output.
\end{proof}
\lean{Graph/OneWay}{88}{92}{Linear.OneWayReplacement.insert_not_input}

\begin{theorem}[Extraction of a filling-local input path]
If $\iota_H(v)$ is operation-reachable in $P$, then $v$ is
operation-reachable inside $H$.
\end{theorem}
\begin{proof}
Choose a reverse operation trace from $\iota_H(v)$ to a global input and
induct on its length, generalizing $v$. If $v$ is already an input of $H$,
the empty local trace is the required witness. Otherwise, in the zero-length
case, $\iota_H(v)$ would be a global input, contradicting the previous lemma.

In the nonempty case, target coverage supplies an edge $f$ of $H$ targeting
$v$. Its inserted copy and the first edge of the global trace target the
same vertex. Unique incoming incidence identifies these two edges. The
first edge is consequently a filling edge and retains its operation label.
Its chosen source is the image of some source $w$ of $f$. Apply the induction
hypothesis to the shorter global trace from that image; this yields a local
operation trace from $w$ to an input of $H$. Prepending $f$ yields the desired
local trace from $v$.

The induction stops on first reaching a filling input. It neither reflects
an arbitrary whole edge path through $\iota_H$ nor assumes injectivity of
that map. This is why the proof remains valid for self-looping holes, unlike
unrestricted path reflection discussed earlier.
\end{proof}
\lean{Graph/OneWay}{98}{124}{Linear.OneWayReplacement.extract}

\Needspace{18\baselineskip}
\subsection{Reflecting host reachability through a one-way filling}
\begin{theorem}[Host term reachability reflects]
Suppose $H$ is one-way and $v$ is a term vertex of $C$. If $\iota_C(v)$ is
operation-reachable in $C[H]$, then $v$ is operation-reachable in $C$.
\end{theorem}
\begin{proof}
Induct on the edge list of a reverse operation trace from $\iota_C(v)$ to
a global input. In the empty case injectivity of the retained-host map
identifies $v$ with that same input.

For a nonempty trace, inspect its first edge. If it is a surviving host
edge, reflect its target incidence by injectivity of $\iota_C$ and express
its source as $\iota_C(w)$. The corresponding host edge is an operation,
so $w$ has term type. Apply the induction hypothesis to the shorter trace
and prepend the host edge.

If the first edge is inserted from $H$, its target is some
$\iota_H(x)=\iota_C(v)$. A filling target cannot be an input. The boundary
classification of replacement therefore forces $x$ to be an output of $H$.
Both replacement maps preserve vertex labels, so this output has term type.
Extraction applied to the original trace now makes it operation-reachable
inside $H$, contradicting one-wayness. Thus only the retained-host case
can occur.
\end{proof}
\lean{Graph/OneWay}{128}{158}{Linear.OneWayReplacement.reflect}

\begin{theorem}[Incoming host paths consist of operations]
Under the same hypotheses, every edge of any host path ending at $v$ is an
operation. The given host path is allowed to contain the replaced edge.
\end{theorem}
\begin{proof}
Reflect operation reachability to $C$, then apply the earlier theorem that
all edges of any path ending at an operation-reachable vertex are operations.
That theorem uses the existing nonempty edge-path representation and allows
repeated edges. No assertion that the supplied host path itself survives
filling is needed.
\end{proof}
\lean{Graph/OneWay}{162}{168}{Linear.OneWayReplacement.path_operations}

\begin{theorem}[One-way filling closure]
If both $C$ and $H$ are one-way, then $C[H]$ is one-way.
\end{theorem}
\begin{proof}
Suppose a term output of $C[H]$ were operation-reachable. It is the retained
copy of the corresponding term output of $C$. Reflection would make that
output operation-reachable in $C$, contradicting the host's one-wayness.
\end{proof}
\lean{Graph/OneWay}{171}{176}{Linear.OneWayReplacement.oneWay_plug}

\Needspace{18\baselineskip}
\subsection{The concrete one-way obstruction: Lemma 26(2)}
\label{sec:one-way-obstruction}
For a focussed host $C:\star\to\epsilon$, replace an edge other than its
focus by a globally focus-free filling $H$. Once the resulting graph is
proved globally hole-free, package it with its retained focus as an actual
state. The constructor requires these focus-freedom and hole-freedom proofs;
it does not infer them merely from the interfaces. Its focus source is
$\iota_C(s)$, where $s$ is the host's focus source.
\lean{Graph/OneWayObstruction}{12}{16}{Linear.Focussed.filledState}

\begin{theorem}[Concrete Lemma 26(2)]
Suppose the focussed host has a path ending at its focus source and that
path contains a non-operation edge. If the filling $H$ is one-way, the
resulting actual state is not rooted.
\end{theorem}
\begin{proof}
Assume the filled state is rooted. The concrete Lemma~23 gives operation
reachability of its focus source $\iota_C(s)$ in the actual filled graph.
The source $s$ of the host focus has term type. Apply the preceding host-path
theorem: every edge of the supplied host path must be an operation. This
contradicts its specified non-operation edge.
\end{proof}
\lean{Graph/OneWayObstruction}{21}{32}{Linear.Focussed.oneWay_not_rooted}

This proves the single-hole one-way obstruction for the current concrete
search semantics. It follows the uniqueness-and-extraction argument of
Proof~C.8, packaged as a reflection lemma rather than a split on whether the
supplied path contains the hole. In fact the Lean result permits replacement
at any non-focus edge with the right interface, so it includes the paper's
hole-labelled context case. The example below uses a genuine hole label.

\begin{theorem}[Paired single-hole obstruction]
If $H$ and $K$ are two admissible fillings of the same focussed host and at
least one is one-way, then at least one of the two filled states is not rooted,
provided the host has the same non-operation path to its focus source.
\end{theorem}
\begin{proof}
Split the hypothesis that $H$ or $K$ is one-way. Apply the preceding theorem
to that filling and introduce the corresponding side of the conclusion.
\end{proof}
\lean{Graph/OneWayObstruction}{37}{50}{Linear.Focussed.oneWay_pair_not_rooted}
This is a single-hole instance of the output-closure argument.
Section~\ref{sec:shallow-specimens} defines concrete pre-templates,
output-closure, and shallow multi-hole specimens and proves their
Lemma~26(3) obstruction. Deep specimen substitution remains open.

\Needspace{18\baselineskip}
\subsection{One-way examples and a filled-state obstruction}
The variable corolla is one-way. Its term output is targeted by a variable
edge. If that output were operation-reachable, the incoming-edge theorem
would force this variable edge to be an operation, a contradiction. By
contrast, the existing operation corolla with two term arguments is not
one-way: its sole edge is a one-step operation path from the input to the
first term output.
\lean{Graph/OneWayExamples}{9}{17}{Linear.Examples.variable_oneWay}
\lean{Graph/OneWayExamples}{21}{24}{Linear.Examples.operation_not_oneWay}

Define an operation signature with no term arguments and one thunk argument,
and take its one-edge corolla. This graph has a term input and only a thunk
output. It is one-way because no output has term type, yet its sole operation
edge explicitly witnesses reachability of its thunk output. Together these
examples verify both the nontriviality of the predicate and the importance
of retaining the paper's term-output restriction.
\lean{Graph/OneWayExamples}{27}{27}{Linear.Examples.thunkOnlyOperations}
\lean{Graph/OneWayExamples}{29}{31}{Linear.Examples.thunkOnlyGraph}
\lean{Graph/OneWayExamples}{34}{39}{Linear.Examples.thunk_only_oneWay}
\lean{Graph/OneWayExamples}{42}{43}{Linear.Examples.thunk_output_reachable}

For a nonvacuous state-level application, start with the earlier closed chain
whose interior focus follows a variable and precedes a value. Relabel that
variable by a hole with one term input and one term output, preserving all
incidences. The other labels remain the focus and the value operation. Their
signatures are unchanged, so the relabelled graph is well formed. Retain the
same distinguished focus, mode, and endpoints. The only other edges are a
hole and an operation, both focus-free; the focus therefore remains globally
unique and satisfies its original endpoint conditions.
\lean{Graph/OneWayExamples}{47}{58}{Linear.Examples.oneWayContext}
\lean{Graph/OneWayExamples}{61}{77}{Linear.Examples.oneWayContext_focus}

Fill the genuine hole by the variable corolla. The exact hole inventory of
shallow substitution shows the resulting graph is hole-free: the filling is
hole-free and no hole survives in the host. Indeed the only host edge that
could carry a hole is exactly the removed edge. Global focus freedom of the
variable filling was already checked.
\lean{Graph/OneWayExamples}{80}{97}{Linear.Examples.oneWayContext_filled_closed}

Finally, take the host path consisting only of its hole edge. Its target is
the focus source, and its label is not an operation. The variable filling is
one-way, so the concrete Lemma~26(2) theorem proves that this actual filled
state is not rooted, for every operation classification. The offending host
edge is removed by substitution: this exercises the new extraction argument,
not merely preservation of a path of retained host edges.
\lean{Graph/OneWayExamples}{101}{109}{Linear.Examples.oneWay_filled_not_rooted}
